\section{Code-Generation Targets}
\label{sec:RustTarget}

The scaffolder is target-parameterised: from a single BPMN model it can render a
generated project in more than one implementation language. The mature target is
Java (Quarkus / SmallRye Reactive Messaging); a Rust target is in progress, and
the experimental CharlotteOS target emits portable Rust procedure contracts plus
cluster deployment descriptors. The Java and ordinary Rust targets emit workers
that speak the same Kafka wire protocol, so a single process
can, in principle, be realised as a mix of Java and Rust workers, and the same
Java monitor observes either.

\subsection{Template organisation}

Template groups are separated by target under \filepath{durga-tools/src/main/resources/}:

\begin{itemize}
\item \filepath{templates-java/} --- the Java target (\filepath{scaffold.stg} and
  companions). This is the default and the only fully implemented target.
\item \filepath{templates-rust/} --- the Rust target (in progress).
\item The CharlotteOS target uses a deterministic descriptor writer rather than
  source templates. It emits a small \texttt{no\_std} Rust contract crate and the
  three reviewable YAML resources under \filepath{charlotte/}.
\end{itemize}

The attribute model passed from the generator (process id, task id, plugin
reference, plugin configuration, channels, lineage) is identical across targets;
only the rendered syntax differs. Loading the template group is the single
switch point in \texttt{BpmnScaffolder}.

\subsection{Shared source of truth}

To keep the targets in lock-step, the following remain single sources of truth
consumed by both:

\begin{itemize}
\item \textbf{Plugin catalog} --- the YAML descriptors under \filepath{plugins/}
  define plugin identity, category, and the typed configuration schema. A model
  scaffolds identically for either target; the same plugin configuration string
  (e.g.\ \texttt{id:order\_id, customer.email:customer\_email}) is honoured by the
  Java and Rust plugin implementations alike.
\item \textbf{Lifecycle wire format} --- the \texttt{ProcessEvent} JSON (camelCase
  fields, SCREAMING\_SNAKE event types and statuses) is byte-compatible across
  targets, so the monitoring topology, alarm model, and dashboard treat
  Rust-generated tasks exactly as Java tasks.
\item \textbf{Chaining topics} --- both targets use identical topic names: a task
  consumes its predecessor's \texttt{\%processId\%\_\%predecessor\%\_output}, a routing
  gateway writes to each target's \texttt{\%processId\%\_\%target\%\_input}, and the
  first task consumes \texttt{\%processId\%\_start}. A Java and a Rust realisation of the
  same model are therefore wire-interchangeable on every topic --- a mixed fleet works,
  and a shadow of either target can validate a production of either target.
\end{itemize}

\subsection{The Rust plugin runtime (\texttt{durga-rust})}

The \filepath{durga-rust/} Cargo crate is the Rust counterpart to the Java
\filepath{durga-runtime} plugin layer. Generated Rust workers depend on it
directly --- there are no adapters onto the Java plugins; each plugin is a native
Rust implementation.

The \texttt{Plugin} trait mirrors the Java \texttt{Plugin} interface, including its
dispatch defaults:

\begin{itemize}
\item \texttt{execute\_bytes(payload, config)} --- binary entry point; the default
  UTF-8 decodes and delegates to \texttt{execute\_str}, then re-encodes. Binary
  plugins override this directly.
\item \texttt{execute\_str(payload, config)} --- text / JSON entry point; text
  plugins override this. The default is unsupported, matching the Java default.
\item \texttt{execute\_with\_result(payload, config, context)} --- the structured
  entry point the generated worker calls, carrying a \texttt{PluginExecutionContext};
  the default wraps \texttt{execute\_bytes} with a content-based idempotency key and the
  \texttt{PAYLOAD} disposition, ignoring the context. Side-effecting plugins override it
  to suppress effects when \texttt{context.validation\_mode()} is true.
\item \texttt{idempotency\_key(payload, config)} --- SHA-256 of payload then
  config, byte-identical to the Java default.
\end{itemize}

\texttt{PluginResult} carries the same \texttt{OutputDisposition}
(\texttt{PAYLOAD}, \texttt{PASSTHROUGH}, \texttt{SIDE\_EFFECT}) and
\texttt{ErrorStrategy} (\texttt{FAIL}, \texttt{SKIP}, \texttt{DLQ}) semantics as
the Java record, so routing and inspection plugins declare \texttt{PASSTHROUGH}
and the generated worker preserves the payload identically on either target.

\subsection{Validation shadows}

Validation mode (Section~\ref{sec:ValidationMode}) is emitted on both targets. With
\texttt{--validation} the scaffolder generates the \emph{whole} process again, wired to
read the live production topics through a dedicated consumer group, suppress side effects
(via the validation \texttt{PluginExecutionContext}), divert each task's output to a
\texttt{-validation} topic, and emit lifecycle events to the per-process
\texttt{process-events-\%processId\%-validation} topic. There is no separate shadow-worker class or validation
binary and no separate candidate wire record: both targets emit ordinary
\texttt{ProcessEvent} records, and because their chaining topics are identical, a shadow
of either target validates a production of either target.

\subsection{CharlotteOS process bundles}
\label{sec:CharlotteTarget}

The experimental \texttt{-{}-target charlotte} path treats generated output as a
process bundle rather than a container image. It preserves the original BPMN
model and its SHA-256 digest, emits one portable procedure module per activity
or exclusive gateway, and records the semantic graph in
\filepath{charlotte/bundle.yaml}. The companion
\filepath{charlotte/deployment.yaml} declares CLS2 artifact identity, singleton
\texttt{PlacementPolicy} fields, generation fencing, and the placeholders that a
trusted build/admission pipeline must replace with artifact and provenance
digests. \filepath{charlotte/capabilities.yaml} is the least-authority review
plan; it is not a bag of credentials.

A Durga procedure is deliberately not a Kafka worker on CharlotteOS. A generic
transactional-step service owns the input delivery and Kafka transaction, calls
the generated procedure endpoint, validates its returned route/output records,
produces to the launch-authorised output routes, includes the input offset, and
then commits. Consequently generated business code cannot leak a delivery,
forget an abort, or select an undeclared topic. Producer and consumer endpoint
roles remain available for applications whose operations are independent, but
a transaction never spans two Kafka service instances.

The low-level Charlotte Kafka profile supports a bounded allow-list of broker
destinations, one fixed consume route, and a bounded, ordered allow-list of
produce topic/partition routes. Producer sequence numbers and transaction
partition state are maintained per route. A version-6 profile is serialized as
one SHA-256-protected object and passed through a kernel-enforced read-only
launch capability; the hash detects corruption while launcher authority
establishes provenance. It fixes the connector identity plus bounded authority
endpoint names and rights ceilings, so generated steps use an independently
named transactional capability rather than a global \texttt{kafka}
publication. Generated descriptors select maxima of 32 broker
endpoints and 64 produce routes; operators may lower those admission ceilings,
while Charlotte rejects declared or actual counts above its hard maxima.
Connector-only SCRAM and mTLS material uses bounded, length-delimited sections;
unknown critical authentication sections fail closed while future optional
sections can be skipped by an older connector.
CharlotteOS now implements the generic \texttt{kafka\_step} runner and its
bounded procedure request/reply ABI. It validates outputs before starting a
transaction, combines admitted records with the input offset, retries transient
procedure failures and timeouts, and transactionally sends terminal, malformed,
or exhausted records to the configured DLQ. Generated deployments still say
\texttt{deployable: false} until activity handlers and controller-supplied
capability bindings exist; the target does not silently weaken Durga's delivery
semantics.

The generated plan treats the broker connector and transactional-step runner
as separate deployment authorities. Only the connector service receives the
profile containing broker addresses, TLS material, and SASL/SCRAM or mTLS
credentials. The version-6 profile publishes a separately named transactional
access point whose consume, produce, and transaction rights are enforced for
every opcode. The runner receives only that access-point capability, and the
activity procedure receives only its bounded invocation.
Credential rotation therefore does not require rebuilding business logic.

Data objects, stores, and associations are retained in the bundle. An S3 data
store becomes an explicit profile request, with credentials referenced from a
launch-time secret profile rather than copied into the generated source. Kafka
transactions do not cover S3 or other external effects, so such activities must
declare an outbox, idempotency, or compensation policy before production use.

\subsection{Status}

Implemented so far: the crate contracts (\texttt{Plugin}, \texttt{PluginResult},
\texttt{PluginExecutionContext}, \texttt{ProcessEvent}, dot-path field helpers), the
pipeline plugin ports, the pure worker decision logic (\texttt{plan\_worker\_output}),
the Rust worker / gateway and project (\texttt{Cargo.toml}, \texttt{lib.rs},
\texttt{build.rs}) templates, and the \texttt{-{}-target rust} switch in the CLI. Rust
workers use a mature Kafka client (\texttt{rdkafka}) and an async runtime
(\texttt{tokio}) for the same consume--transform--emit loop the plain-Kafka Java
worker performs. Rust remains the less-mature target: it currently covers plugin
tasks, custom-task stubs (with contract modules and \texttt{build.rs} entries),
exclusive gateways, and validation mode in a start/tasks/gateways/end topology, and
reports and skips unsupported constructs.

The CharlotteOS target currently implements bundle, deployment, capability, and
portable procedure-contract generation. It does not yet emit runnable ELFs or
submit a cluster deployment. This boundary is intentional: signing, provenance,
capability grants, placement realisation, and transactional-step supervision are
trusted deployment stages rather than code-generator side effects.
